\documentclass[runningheads,orivec]{llncs}
\usepackage[T1]{fontenc}
\usepackage{graphicx}
\usepackage{booktabs}
\usepackage{multirow}
\usepackage{url}
\usepackage{appendix}
\usepackage{subcaption}
\usepackage{float}

\input{ECML_PKDD_2025_Author_Kit/metadata}

\newcommand{\rankingagent}{\textit{Ranking agent}}
\newcommand{\technicalagent}{\textit{Technical agent}}
\newcommand{\technicalagents}{\textit{Technical agents}}
\newcommand{\locationagent}{\textit{Location agent}}
\newcommand{\energyagent}{\textit{Energy agent}}
\newcommand{\proximityagent}{\textit{Proximity agent}}
\newcommand{\regulatoryagent}{\textit{Regulatory agent}}
\newcommand{\buildingagent}{\textit{Building agent}}
\newcommand{\sqlagent}{\textit{SQL agent}}
\newcommand{\evaluationagent}{\textit{Evaluation agent}}
\newcommand{\brokeragent}{\textit{Broker agent}}
\newcommand{\filteringmode}{\textit{Filtering mode}}
\newcommand{\rankingmode}{\textit{Ranking mode}}
\newcommand{\gptoss}{\texttt{gpt-oss-120b}}
\newcommand{\gemma}{\texttt{gemma3-27b}}
\newcommand{\qwen}{\texttt{qwen3-8b}}
\newcommand{\gptfive}{\texttt{gpt-5.4}}
\newcommand{\combqueries}{$\mathcal{Q}_\text{comb}$}
\newcommand{\completequeries}{$\mathcal{Q}_\text{full}$}

\usepackage{amssymb}
\usepackage{amsmath}
\usepackage{cleveref}
\usepackage{listings}
\usepackage{tabularx}
%%%%%%%%%%%%%%%%%%%%%

\usepackage{xcolor}
\usepackage{ifthen}
\usepackage[misc]{ifsym}

\usepackage[style=numeric, citestyle=numeric-comp, sorting=none, backend=bibtex, maxcitenames=2, mincitenames=1]{biblatex}



\begin{document}

\title{\paperTitle: Appendix}

\titlerunning{\methodACR\ - Appendix}

\author{\paperAuthor}
\authorrunning{\paperAuthorRunning}

\institute{\paperAffiliation}

\maketitle

\centerline{\textbf{Supplementary Material}}

% %%%%%%%%%%%%%%%%%%%%%%%%%%%%%%%%%%%%%%%%%%%%%%
\appendix

\section{Visual inspection}

\begin{figure}[H]
    \centering
    
    \begin{subfigure}{0.48\linewidth}
        \centering
        \includegraphics[width=\linewidth]{ECML_PKDD_2025_Author_Kit/Supplementary/images/app/login.png}
        \caption{Login page.}
        \label{fig:login}
    \end{subfigure}
    \hfill
    \begin{subfigure}{0.48\linewidth}
        \centering
        \includegraphics[width=\linewidth]{ECML_PKDD_2025_Author_Kit/Supplementary/images/app/home.png}
        \caption{Home page.}
        \label{fig:home}
    \end{subfigure}

    \vspace{0.5em}

    \begin{subfigure}{0.48\linewidth}
        \centering
        \includegraphics[width=\linewidth]{ECML_PKDD_2025_Author_Kit/Supplementary/images/app/map.png}
        \caption{Map visualization.}
        \label{fig:map}
    \end{subfigure}
    \hfill
    \begin{subfigure}{0.48\linewidth}
        \centering
        \includegraphics[width=\linewidth]{ECML_PKDD_2025_Author_Kit/Supplementary/images/app/inspector.png}
        \caption{Inspector view.}
        \label{fig:inspector}
    \end{subfigure}
    
    % \vspace{0.5em}

    % \begin{subfigure}{\linewidth}
    %     \centering
    %     \includegraphics[width=\linewidth]{ECML_PKDD_2025_Author_Kit/images/app/details.png}
    %     \caption{Property details view.}
    %     \label{fig:details}
    % \end{subfigure}

    \caption{User interface of the decision-support system: authentication, query interface, spatial visualization, asset inspection, and agent-level reasoning transparency. Inspector view shows agent reasoning and intermediate results.}
    \label{fig:app_overview}
\end{figure}


%%%%%%%%%%%%%%%%%%%%%%%%%%%%%%%%%%%%%%%%%%%%%%
\section{Extended results}
%%%%%%%%%%%%%%%%%%%%%%%%%%%%%%%%%%%%%%%%%%%%%%
\subsection{Combinatorial query generation protocol}

To systematically evaluate the pipeline, we constructed a controlled set of synthetic user queries via combinatorial variation of six input components: typology (\(I\)), intended use (\(D\)), point of interest (\(L\)), building constraints (\(C_b\)), energy constraints (\(C_e\)), and ancillary services (\(C_p\)). The full configuration space yields 486 queries (\(\mathcal{Q}_{\mathrm{comb}}\)), spanning both mandatory and optional components (Table~\ref{tab:query-components}).

Optional fields were independently activated to simulate realistic variability in user requests, producing queries of heterogeneous structural complexity. Among these, 64 queries (\(\mathcal{Q}_{\mathrm{full}}\)) simultaneously activate all five technical agents, representing the most constrained and computationally demanding scenarios. This controlled design enables a systematic evaluation of both routing robustness and multi-agent interaction as query complexity increases.


\begin{table}[H]
\caption{Input components and their values for the combinatorial generation of the evaluation queries. Reported dimensions are typology $I$, intended use $D$, point of interest $L$, building constraints $C_b$, energy constraints $C_e$, and ancillary services $C_p$. Optional fields are marked with *.}
\label{tab:query-components}
\centering
\begin{tabular}{ll}
\toprule
\textbf{Component} & \textbf{Options} \\ \midrule
Property typology ($I$) & building; dwelling \\ 
Intended use ($D$)* & co-housing; elderly care center \\ 
Point of interest ($L$)* & Landmark 1; Landmark 2 \\ 
% Point of interest ($L$) & Palazzo Nuovo; Porta Susa; none \\ 
Surface area ($C_b$)* & between 50 and 80 m\(^2\); over 1000 m\(^2\) \\ 
Energy class ($C_e$)* & F or G; C or better \\ 
Ancillary services ($C_p$)* & public transport; gyms and sports centers \\ \bottomrule
\end{tabular}
\end{table}

\subsection{Extended ablation study}

% \subsubsection{Agent sensitivity, ranking, and knowledge contribution.}
Table \ref{tab:ablation_summary} shows a leave-one-out ablation over the \completequeries\ to quantify each component's contribution.
Ranking IoU is computed between the top-10 of the ablated configuration and the full pipeline; a lower value indicates a higher discriminative contribution.
As shown in \Cref{tab:ablation_summary}, for \gptoss\ the \energyagent\ is the most impactful component (IoU $= 0.082$), followed by the \proximityagent\ and \locationagent; for \gemma\ the \buildingagent\ drives the largest perturbation (IoU $= 0.072$); for \qwen\ the \locationagent\ shows the highest sensitivity (IoU $= 0.145$), followed by the \proximityagent\ (IoU $= 0.154$).
Removing the \rankingagent's weighting (uniform weights) does not alter the retrieved top-10 (IoU $= 1.000$ for all models). This result is consistent with the structure of the evaluation dataset: the \combqueries\ are generated from a small combinatorial grid of query templates, which produces a relatively uniform distribution of agent activations. Under such conditions, harmonic weights and uniform weights yield similar score distributions, and the top-10 ranking is not sufficiently sensitive to weight perturbations to produce measurable divergence. On queries with a stronger imbalance across agent importance, the \rankingagent's semantic prioritisation would be expected to produce a more pronounced effect.
Similarly, withholding data statistics (\textit{no knowledge}) leaves the retrieved set unchanged (IoU $= 1.000$), but collapses the constraint tightness from $0.279$ to $0.000$. This suggests that on the structured EPC dataset, the LLM already possesses sufficient implicit knowledge of plausible numerical ranges to generate syntactically valid SQL filters even without explicit statistical grounding; however, without injected quartiles, the generated thresholds lack discriminative calibration to the actual data distribution, producing vacuous or maximally permissive constraints.

\begin{table}[h]
\centering
\caption{Ranking IoU between each ablated configuration and the full pipeline, computed on the \completequeries.}
\label{tab:ablation_summary}
\begin{tabular}{lccc}
\toprule
\textbf{Configuration} & \textbf{\gptoss} & \textbf{\gemma} & \textbf{\qwen} \\
\midrule
No \energyagent           & 0.082 & 0.175 & 0.311 \\
No \proximityagent        & 0.161 & 0.214 & 0.154 \\
No \locationagent         & 0.245 & 0.169 & 0.145 \\
No \buildingagent         & 0.405 & 0.072 & 0.275 \\
No \regulatoryagent       & 0.482 & 0.692 & 0.561 \\
\bottomrule
\end{tabular}
\end{table}

\section{Data aggregation policy}
%%%%%%%%%%%%%%%%%%%%%%%%%%%%%%%%%%%%%%%%%%%%%%

This appendix provides a complete description of the data
processing pipeline used to ingest, parse, enrich, and consolidate
Italian Energy Performance Certificates (EPC) for the municipality of Turin (cadastral
code \texttt{L219}). The pipeline
operates on a corpus of XML files conforming to the CSI/SICEE 2015
schema. It produces one record
per certified real-estate unit, augmented with geographic,
urban-amenity, and energy-quality scores.
The pipeline comprises five sequential stages, summarised in
Table~\ref{tab:stages}.

\begin{table}[h]
\centering
\caption{Summary of pipeline stages.}
\label{tab:stages}
\begin{tabularx}{\textwidth}{clX}
\toprule
\textbf{Stage} & \textbf{Name} & \textbf{Description} \\
\midrule
1 & XML Parsing         & Extraction of certificate fields via XPath from SICEE-schema XML files \\
2 & Geocoding           & Resolution of postal addresses to WGS-84 coordinates using a local OSM address database \\
3 & Geographic enrichment & Assignment of OMI real-estate zones and computation of urban-amenity proximity scores \\
4 & Energy scoring      & Derivation of composite energy-quality indicators \\
5 & Deduplication and aggregation & Cadastral deduplication and synthesis of building-level meta-records \\
\bottomrule
\end{tabularx}
\end{table}

% -------------------------------------------------------
\subsection{Stage 1: XML parsing}
\label{sec:parsing}
% -------------------------------------------------------

Each EPC file is an XML document governed by the namespace \\
\texttt{http://www.csi.it/sicee/siceeweb/xml/xmlapecompleto2015/data}. Six \\groups of attributes are extracted via XPath expressions.

\paragraph{Identity and location.}
Street address, civic number, municipality, climatic zone, construction
year, floor level, and GIS coordinates (latitude/longitude) where
embedded in the certificate.

\paragraph{Global energy performance.}
Energy class (scale A4--G), non-renewable global energy performance
index ($EP_{gl,nren}$, kWh/m\textsuperscript{2}$\cdot$year),
renewable component \\ ($EP_{gl,ren}$), and CO\textsubscript{2} emissions.

\paragraph{Energy carriers.}
Fourteen energy vectors are read (Table~\ref{tab:carriers}). For
fossil fuels and biomass, raw consumption values (expressed in
volume or mass units) are converted to kWh equivalents by
multiplication with lower heating value factors ($H_i$, also listed
in Table~\ref{tab:carriers}). The carrier with the largest converted
consumption is designated the \emph{primary energy carrier}.

\begin{table}[h]
\centering
\caption{Energy carriers and lower heating value conversion factors.
Entries marked ``---'' are already expressed in kWh.}
\label{tab:carriers}
\begin{tabular}{llc}
\toprule
\textbf{Carrier} & \textbf{Unit (XML)} & $H_i$ \textbf{(kWh/unit)} \\
\midrule
Grid electricity      & kWh   & --- \\
Natural gas           & m\textsuperscript{3} & 9.94 \\
LPG                   & kg    & 12.778 \\
Coal                  & kg    & 7.917 \\
Diesel                & kg    & 11.87 \\
Fuel oil              & kg    & 11.75 \\
Solid biomass         & kg    & 4.67  \\
Liquid biomass        & kg    & 7.50  \\
Gaseous biomass       & m\textsuperscript{3} & 6.40  \\
Solar photovoltaic    & kWh   & --- \\
Solar thermal         & kWh   & --- \\
Wind                  & kWh   & --- \\
District heating      & kWh   & --- \\
District cooling      & kWh   & --- \\
\bottomrule
\end{tabular}
\end{table}

\paragraph{Envelope quality.}
The XML encodes winter and summer envelope performance as integer
codes (0, 1, 2), mapped to qualitative ratings: \textit{satisfactory},
\textit{neutral}, and \textit{poor}.

\paragraph{Building services and heating systems.}
Boolean flags record the presence of heating, domestic hot water
(DHW), cooling, mechanical ventilation, lighting, and
lift/transportation services. For heating and DHW, the technology
code and installation year of the primary generator are extracted
and mapped to human-readable descriptions via a look-up table
(e.g.\ codes \texttt{4}--\texttt{15} $\to$ heat pump, code
\texttt{1} $\to$ condensing boiler).

\paragraph{Cadastral data.}
Municipality code, cadastral sheet, parcel, and sub-unit
are extracted. When the XML records a range
of sub-units (\texttt{subDA}--\texttt{subA}), both bounds are
retained as a comma-separated string.

\paragraph{Improvement recommendations.}
Target energy class and estimated payback period from the
certificate's recommendation section.

% -------------------------------------------------------
\subsection{Stage 2: local OSM geocoding}
\label{sec:geocoding}
% -------------------------------------------------------

EPC certificates frequently omit GIS coordinates or contain
erroneous ones. A local geocoding service based on OpenStreetMap
data is therefore used to derive or improve coordinate assignments
without reliance on external APIs.
A local address database is built from OSM data covering
north-west Italy, retaining all nodes that carry street name,
house number, and city tags, indexed by street name for efficient lookup.

\subsubsection{Address normalisation}

Raw address strings from the EPC certificates are normalised
through the following steps, applied in sequence:

\begin{enumerate}
  \item Removal of leading postal codes (five-digit sequences).
  \item Removal of city-name tokens (e.g.~\textit{Turin}, abbreviated \textit{TO})
        at the start or end of the string.
  \item Removal of stray leading numeric tokens.
  \item Expansion of abbreviated street-type prefixes
        (e.g.\ \texttt{C.SO} $\to$ \texttt{CORSO} [avenue],
        \texttt{P.ZZA} $\to$ \texttt{PIAZZA} [square]).
  \item Extraction of civic numbers embedded in the street string
        (e.g.\ ``\texttt{CORSO TRAPANI, 133}'' $\to$
        street = \texttt{corso trapani}, number = \texttt{133}).
\end{enumerate}

\subsubsection{Candidate retrieval and scoring}

Up to 60 candidate records are retrieved for the normalised street.
If no candidates are found, a fallback query drops the street-type prefix
(retaining only the toponym). Each candidate is scored according to:

\begin{equation}
  s = s_\text{city} + s_\text{street} + s_\text{number}
  \label{eq:geo_score}
\end{equation}

where the component terms are defined in
Table~\ref{tab:geo_scoring}.

\begin{table}[h]
\centering
\caption{Geocoding scoring components.}
\label{tab:geo_scoring}
\begin{tabular}{lll}
\toprule
\textbf{Component} & \textbf{Condition} & \textbf{Points} \\
\midrule
\multirow{2}{*}{City ($s_\text{city}$)}
  & Exact match & $+100$ \\
  & No match    & $0$    \\
\midrule
\multirow{3}{*}{Street ($s_\text{street}$)}
  & Exact match          & $+150$ \\
  & Substring match      & $+50$  \\
  & No overlap           & $0$    \\
\midrule
\multirow{3}{*}{Number ($s_\text{number}$)}
  & Exact match          & $+300$ \\
  & Prefix match         & $+150$ \\
  & Mismatch             & $-100$ \\
\bottomrule
\end{tabular}
\end{table}

The candidate with the highest score is selected. Its coordinates
replace the XML-supplied values if (a) the score exceeds a quality
threshold of 450 points, or (b) the XML contained no valid
coordinates.

% -------------------------------------------------------
\subsection{Stage 3: geographic enrichment}
\label{sec:enrichment}
% -------------------------------------------------------

The Italian Revenue Agency publishes OMI (``Osservatorio del Mercato Immobiliare'') zones, i.e. cadastral micro-zones used for official real-estate valuations.
For each property, a point-in-polygon test against
OMI zone boundaries assigns the corresponding zone code.

\subsubsection{Urban amenity proximity scoring}

\paragraph{POI dataset}

Points of Interest (POIs) are extracted from OpenStreetMap (OSM) data and organised
into six thematic categories: \textit{health}, \textit{mobility},
\textit{greenery}, \textit{sport}, \textit{commerce},
and \textit{education}.
The extraction logic processes the three primary OSM geometry types: nodes, ways, and
relations. Nodes are handled directly as point geometries; the spatial coordinates of
ways and relations are represented by their centroids, computed by averaging the
coordinates of their constituent nodes.

A tag-based categorisation mechanism maps OSM tags (e.g., \texttt{amenity},
\texttt{shop}, \texttt{leisure}) to the six target categories, covering over 230
distinct tags. For complex geometries (ways and relations), a consistency check is
performed across member tags: if all members share a common category, it is assigned
to the aggregate entity; otherwise the category is inherited from the top-level tags
of the way or relation. A deduplication step ensures that nodes already accounted for
within ways or relations are not re-processed as independent POIs. Finally, specific
refinement rules resolve semantic overlaps (e.g., swimming pools are reallocated
across categories depending on analysis context).

Within each category, individual amenity types are classified by expected service
range into three distance-threshold classes:
\textit{short} ($S = 340\,\text{m}$, $\approx$5-min walk),
\textit{medium} ($S = 680\,\text{m}$, $\approx$10 min), and
\textit{long} ($S = 1{,}360\,\text{m}$, $\approx$20 min),
adhering to the ``15-minute city'' planning paradigm.

\paragraph{Score computation}

All POI coordinates are stored in a spatial index (haversine metric).
For a property located at coordinates $(\varphi, \lambda)$, all
POIs within a maximum radius of 5\,km are retrieved. The
contribution of a single POI at distance $x$ (km) is modelled by
a cubic Gaussian decay:

\begin{equation}
  f(x) = \exp\!\left[-\left(\frac{x}{S \cdot U}\right)^{\!3}\right]
  \label{eq:decay}
\end{equation}

where $S$ is the base service radius for the amenity's threshold
class and $U$ is an urban-density multiplier depending on the
urban zone type of the property (Table~\ref{tab:urban_mult}).

\begin{table}[h]
\centering
\caption{Urban density multipliers by urban zone type.}
\label{tab:urban_mult}
\begin{tabular}{cc}
\toprule
\textbf{Zone} & $U$ \\
\midrule
B & 1.0 \\
C & 1.5 \\
D & 2.5 \\
E & 4.0 \\
\bottomrule
\end{tabular}
\end{table}

Individual POI scores are accumulated by amenity type. The
category score $c_k$ is the mean across all amenity types
$\{a_j\}$ belonging to category $k$:

\begin{equation}
  c_k = \frac{1}{|A_k|}\sum_{j \in A_k}
        \sum_{\substack{i\,:\,\text{cat}(i)=k \\ \text{am}(i)=j}}
        f(x_i)
  \label{eq:cat_score}
\end{equation}

The six raw category scores are percentile-normalised to the
range $[0, 100]$ across the full dataset.

The choice of a cubic exponent ($k=3$) in Eq.~\ref{eq:decay}, rather than a
standard Gaussian ($k=2$) or linear ($k=1$) decay, is deliberate. It creates a
\emph{plateau of indifference} near the origin where $f(x)\approx 1$, reflecting
human perception: services within an immediate walkable radius tend to be perceived
as equally convenient regardless of exact distance. Beyond this plateau, the
higher-order exponent induces a sharper utility drop-off, modelling the threshold
effect whereby an amenity that transitions from walkable to transit-required loses
value more abruptly than a smooth quadratic decay would predict. The urban multiplier
$U$ (Table~\ref{tab:urban_mult}) ranges from $1.0$ in dense, high-walkability central
areas to $4.0$ in suburban environments with higher car dependency, and is calibrated
against density expectations in major Italian metropolitan areas.

% -------------------------------------------------------
\subsection{Stage 4: energy quality scoring}
\label{sec:scoring}
% -------------------------------------------------------

Four ordinal scores are computed for each property from the
parsed EPC fields (Table~\ref{tab:energy_scores}). The composite
score $S_\text{total}$ is their sum (range: 4--20).

\begin{equation}
  S_\text{total} = S_\text{class} + S_\text{system} +
                   S_\text{envelope} + S_\text{renewables}
\end{equation}

\begin{table}[h]
\centering
\caption{Energy quality score mapping.}
\label{tab:energy_scores}
\begin{tabularx}{\textwidth}{lXc}
\toprule
\textbf{Score} & \textbf{Condition} & \textbf{Value} \\
\midrule
\multirow{3}{*}{$S_\text{class}$}
  & Energy class starts with A & 5 \\
  & Energy class is B, C, D, or E & 3 \\
  & Energy class is F or G (or missing) & 1 \\
\midrule
\multirow{3}{*}{$S_\text{system}$}
  & Heating system is heat pump or district heating & 5 \\
  & Heating system is condensing boiler or biomass  & 3 \\
  & Standard boiler or other  & 1 \\
\midrule
\multirow{3}{*}{$S_\text{envelope}$}
  & Envelope quality \textit{satisfactory} & 5 \\
  & Envelope quality \textit{neutral}       & 3 \\
  & Envelope quality \textit{poor} (or missing) & 1 \\
\midrule
\multirow{2}{*}{$S_\text{renewables}$}
  & Solar PV or solar thermal consumption $> 0$ & 5 \\
  & Otherwise & 1 \\
\bottomrule
\end{tabularx}
\end{table}

% -------------------------------------------------------
\subsection{Stage 5: deduplication and meta-property generation}
\label{sec:dedup}
% -------------------------------------------------------

\subsubsection{Cadastral deduplication}

Multiple EPC certificates may exist for the same real-estate unit,
identified by the triple
\textit{sheet}/\textit{parcel}/\textit{sub-unit}.
The deduplication procedure retains, for each unit, only the most
recent certificate (sorted by \texttt{certificate\_date} descending).
An additional address-consistency filter ensures that only records
sharing the address of the most recent entry are considered, to
prevent erroneous merging of data from different units assigned to
the same cadastral triple.

\subsubsection{Meta-property records}

For each (sheet, parcel) pair containing more than one
sub-unit, a synthetic \emph{meta-property} record is generated to
represent the entire building parcel. The aggregation rules are:

\begin{itemize}
  \item \textbf{Floor area}: sum of individual sub-unit areas.
  \item \textbf{Property type}: type with the greatest aggregate
        floor area (surface-weighted mode).
  \item \textbf{Energy class}: class with the greatest aggregate
        floor area (surface-weighted mode).
  \item \textbf{Energy consumption indices}
        ($EP_{gl,nren}$, $EP_{gl,ren}$, CO\textsubscript{2}):
        minimum across sub-units (best-performing unit).
  \item \textbf{Energy sub-scores} ($S_\text{system}$,
        $S_\text{envelope}$, $S_\text{renewables}$): maximum
        across sub-units.
  \item \textbf{All remaining fields}: mode across sub-units.
\end{itemize}

Meta-property records receive identifiers of the form
\texttt{M00001}, \texttt{M00002}, $\ldots$ and the boolean flag
\texttt{is\_meta\_property\,=\,True}. They are appended to the dataset
after all individual unit records.

% -------------------------------------------------------
\subsection{Output schema}
\label{sec:output}
% -------------------------------------------------------

The final dataset contains the columns listed in
Table~\ref{tab:schema}. Columns retained exclusively for intermediate processing
(\texttt{heating\_type}, \texttt{winter\_quality}, \\ \texttt{renewable\_sources},
\texttt{floor}, etc.) are dropped from the final output.
Records lacking an OMI zone assignment are removed.
The six amenity scores are \\ percentile-normalised to
$[0, 100]$ and rounded to two decimal places.

\begin{table}[h]
\centering
\caption{Selected output columns and their descriptions.}
\label{tab:schema}
\begin{tabularx}{\textwidth}{llX}
\toprule
\textbf{Column} & \textbf{Type} & \textbf{Description} \\
\midrule
\texttt{id}                         & string  & Unique record identifier \\
\texttt{is\_meta\_property}          & bool    & True for building-level aggregate records \\
\texttt{num\_units\_per\_parcel}      & int     & Number of sub-units in the parcel \\
\texttt{cadastral\_sheet}            & int     & Cadastral sheet \\
\texttt{cadastral\_parcel}           & int     & Cadastral parcel \\
\texttt{cadastral\_subunit}          & int     & Cadastral sub-unit (0 for meta-properties) \\
\texttt{address}                    & string  & Street address \\
\texttt{latitude}                   & float   & WGS-84 latitude \\
\texttt{longitude}                  & float   & WGS-84 longitude \\
\texttt{omi\_zone}                   & string  & OMI valuation zone code \\
\texttt{floor\_area\_sqm}            & float   & Heated useful floor area (m\textsuperscript{2}) \\
\texttt{property\_type}             & string  & Property use type (DPR~412 classification) \\
\texttt{construction\_epoch}         & string  & Construction year or decade \\
\texttt{certificate\_date}           & date    & Certificate issue date \\
\texttt{epc\_energy\_class}          & string  & Energy class (A4--G) \\
\texttt{epglnren\_epc}               & float   & $EP_{gl,nren}$ (kWh/m\textsuperscript{2}$\cdot$year) \\
\texttt{epglren\_epc}                & float   & $EP_{gl,ren}$ (kWh/m\textsuperscript{2}$\cdot$year) \\
\texttt{co2\_emissions}              & float   & CO\textsubscript{2} emissions (kg/m\textsuperscript{2}$\cdot$year) \\
\texttt{epc\_score\_class}           & int     & Energy class score (1, 3, or 5) \\
\texttt{epc\_score\_system}          & int     & Heating system score (1, 3, or 5) \\
\texttt{epc\_score\_envelope}        & int     & Envelope quality score (1, 3, or 5) \\
\texttt{epc\_score\_renewables}      & int     & Renewables score (1 or 5) \\
\texttt{epc\_score\_total}           & int     & Composite energy score (4--20) \\
\texttt{health}    & float   & Health amenity score (percentile 0--100) \\
\texttt{education} & float   & Education amenity score (percentile 0--100) \\
\texttt{mobility}  & float   & Mobility amenity score (percentile 0--100) \\
\texttt{sport}     & float   & Sport amenity score (percentile 0--100) \\
\texttt{commerce}  & float   & Commerce amenity score (percentile 0--100) \\
\texttt{greenery}  & float   & Greenery amenity score (percentile 0--100) \\
\bottomrule
\end{tabularx}
\end{table}

% POIs description integrated into Stage 3 --- Urban amenity proximity scoring above.

\end{document}
